\section{Conclusion}

\begin{frame}{结论}
  \begin{block}{核心发现}
    数字抵押品并不依赖“物理回收资产”，而是依赖“暂停服务流”来同时缓解\textbf{逆向选择}与\textbf{道德风险}。
  \end{block}

  \begin{itemize}
    \item 接受率下降约 7 个百分点，表明数字抵押带来正向筛选。
    \item 还款率提高约 11 个百分点、全额偿还率提高约 19 个百分点，说明数字抵押显著改善履约与盈利性。
    \item 学费贷款提高了入学率和学校支出，且没有显著恶化家庭资产负债表。
  \end{itemize}

  \scriptsize
  这意味着在弱执行环境中，数字化的使用权控制本身就可能成为一种有效的金融合约基础设施。
\end{frame}

\begin{frame}{福利含义与代价}
  \begin{block}{为什么结果重要}
    结果表明，一部分家庭原本缺乏其他教育融资渠道。样本中约 15\% 的家庭报告在过去 12 个月内曾被拒绝贷款。
  \end{block}

  \begin{block}{但数字抵押并非没有成本}
    \begin{itemize}
      \item 生产与整合锁定技术本身存在固定成本。
      \item 锁定设备会带来事后效率损失: 有担保贷款的中位数家庭在放款后前 200 天中约有 50 天被锁定。
      \item 因此，最优合约不是“锁得越严越好”，而是要在扩大信贷可得性与减少锁定扭曲之间权衡。
    \end{itemize}
  \end{block}

  \scriptsize
  更好的合同设计，可能来自更灵活的宽限期、重组机制或对高延续价值客户的差异化安排。
\end{frame}
